% Generated by tools/edn_latex.py from reports/edn.md.
% Do not edit: change reports/edn.md and run `make edn`.
\ednentry{
  id = {EDN-106},
  title = {The images install their environment from a runtime-only export of the lock},
  date = {2026-11-11},
  milestone = {M5: Model Packaging},
  topic = {Containerization, Dependency Management, CI/CD},
  participants = {Lukas Häußler (decided), AI agents (options, measurements, implementation, review)},
  decision = {A stage of each Dockerfile exports the part of its lock the image installs, \texttt{uv export -\allowbreak{}-\allowbreak{}locked} as a requirements file with hashes, without the project and the groups the image leaves out.
The environment is installed from that file alone with \texttt{uv pip sync -\allowbreak{}-\allowbreak{}require-\allowbreak{}hashes}, so a lock change in a group only keeps the install step cached and the environment layer's digest unchanged.},
  rationale = {\emph{Alternatives considered:}
\begin{itemize}[leftmargin=*, topsep=2pt]
\item A \texttt{pylock.\allowbreak{}toml} export, taken first: installing from it is experimental in uv 0.13.6 and warns on every build.
\item \texttt{uv sync} on \texttt{uv.\allowbreak{}lock}, as before: every weekly dev bump makes the VM pull the environment again.
\item \texttt{uv sync} keyed on the export: impossible, as \texttt{uv sync} reads \texttt{uv.\allowbreak{}lock}, which keys the step on the whole lock.
\end{itemize}
\par
\emph{Why:} An uncached install never writes the same bytes twice, so only a cached step keeps the layer.
Measured with a fresh builder and a cache, a \texttt{types-\allowbreak{}tqdm} bump kept the API's 176 MB layer and a \texttt{types-\allowbreak{}pyyaml} bump the monitoring image's 370 MB one, while the previous Dockerfiles rebuilt both and a \texttt{pydantic} bump still rebuilds.
The environments are unchanged apart from the bytecode.
The cost is that the install finds its files through PyPI's index, each checked against the locked hash.},
  ai = {Alternative generation, Alternative assessment, Recommendation, Solution generation},
  response = {Accepted with modifications},
  assessment = {An AI worker proposed the export, measured layer digests with and without the change and compared the environments; an independent AI review found that the \texttt{pylock.\allowbreak{}toml} install was experimental and that the docs left out base and uv image updates.
Lukas first took the \texttt{pylock.\allowbreak{}toml} recommendation, then, asked again with that fact, the requirements file.},
  aievidence = {Claude Code session of 2026-11-11.
Question (translated from German): ``\#346 once more, with a new fact: the review of PR \#359 showed that uv marks installing from a pylock.toml (\texttt{uv pip sync} with pylock) as experimental in uv 0.13.6 (a warning in every build). That was not in my question earlier, so I ask again. Which way should it be?''
Orchestrator briefing (core, paraphrased): ``Keep the environment layer of both images byte-identical when a lock change touches only groups the image does not install; verify that BuildKit keys the install on a runtime-only export and that the environment stays the same; stop with DECISION NEEDED on an EDN decision.''
Lukas answered: ``requirements.txt export (recommended): the stable uv way with the same effect: export only the runtime, \texttt{uv pip sync -\allowbreak{}-\allowbreak{}require-\allowbreak{}hashes}; the reviewer measured: the export is byte-identical for dev bumps, the same 37 packages, no warning, every file checked by hash. Con: files are resolved through the index (requests to PyPI during the build).'' (translated from German)},
  evidence = {\href{https://github.com/mlops-2627q1-mds-upc/MLOPS_Recommenditos/issues/346}{Issue \#346}, \href{https://github.com/mlops-2627q1-mds-upc/MLOPS_Recommenditos/pull/359}{PR \#359}, \href{https://github.com/mlops-2627q1-mds-upc/MLOPS_Recommenditos/blob/main/docs/docs/deployment.md\#continuous-delivery}{What rebuilds the environment layer}, EDN-47, EDN-105.},
}
